\documentclass[letterpaper, 12pt]{article}
% ============================================================
%  CONFIGURACION COMUN PARA SOLUCIONARIOS PEP 1
% ============================================================

\usepackage[utf8]{inputenc}
\usepackage[spanish]{babel}
\usepackage{amsmath, amssymb, amsfonts}
\usepackage{geometry}
\usepackage{fancyhdr}
\usepackage{enumitem}
\usepackage{graphicx}
\usepackage{hyperref}
\usepackage{needspace}
\usepackage{parskip}
\usepackage{xcolor}
\usepackage{tcolorbox}
\usepackage{eso-pic}
\usepackage{tikz}
\tcbuselibrary{skins, breakable}

\geometry{letterpaper, margin=1.7cm, top=3cm, bottom=2.5cm, headheight=2cm}

\definecolor{celesteSuave}{HTML}{F0F7FF}
\definecolor{azulFuerte}{HTML}{1A5276}
\definecolor{enlace}{HTML}{1A5276}

\newcommand{\R}{\mathbb{R}}
\newcommand{\Dom}{\operatorname{Dom}}
\newcommand{\Rec}{\operatorname{Rec}}
\newcommand{\Cod}{\operatorname{Cod}}

\newcommand{\logoup}{../../../../../template/images/logo_up.png}
\newcommand{\logousach}{../../../../../template/images/logo_usach.png}
\newcommand{\logofondo}{../../../../../template/images/logo_up.png}
\newcommand{\institucion}{Usach Premium}
\newcommand{\curso}{C\'alculo I}
\newcommand{\autor}{Jaime Riquelme}
\newcommand{\semestre}{1er. Semestre de 2026}
\newcommand{\fraseportada}{``Por y para estudiantes''}
\newcommand{\webinstitucion}{www.usachpremium.cl}
\newcommand{\instagraminstitucion}{https://www.instagram.com/usach.premium}
\newcommand{\transparencia}{0.05}

\hypersetup{
    colorlinks=true,
    linkcolor=enlace,
    urlcolor=enlace,
    citecolor=enlace,
    pdfborder={0 0 0},
}

\pagestyle{fancy}
\fancyhf{}
\fancyhead[L]{\includegraphics[height=1.2cm]{\logoup}}
\fancyhead[R]{%
    \begin{minipage}[b]{0.42\textwidth}
        \raggedleft
        \fontsize{9pt}{11pt}\selectfont
        \textbf{\color{azulFuerte}\institucion} \\
        \curso\ -- Solucionario PEP 1 \\
        \textit{\color{gray}@usach.premium}
    \end{minipage}\hspace{0.1cm}%
    \includegraphics[height=1.2cm]{\logousach}%
}
\fancyfoot[L]{\fontsize{9pt}{11pt}\selectfont \textit{\fraseportada}}
\fancyfoot[R]{\fontsize{9pt}{11pt}\selectfont P\'agina \thepage}
\renewcommand{\headrulewidth}{0.4pt}
\renewcommand{\footrulewidth}{0.4pt}

\fancypagestyle{plain}{
    \fancyhf{}
    \renewcommand{\headrulewidth}{0pt}
    \renewcommand{\footrulewidth}{0pt}
}

\newcommand{\MarcaAgua}{
    \AddToShipoutPicture{
        \AtPageCenter{
            \begin{tikzpicture}[remember picture, overlay]
                \node[opacity=\transparencia] at (0,0)
                    {\includegraphics[width=0.7\textwidth]{\logofondo}};
            \end{tikzpicture}
        }
    }
}

\newtcolorbox{solucionbox}[1]{
    colback=white,
    colframe=azulFuerte,
    coltitle=white,
    colbacktitle=azulFuerte,
    title=\textbf{#1},
    title after break=\textbf{#1} (continuaci\'on),
    fonttitle=\bfseries,
    arc=4pt,
    boxrule=1pt,
    enhanced,
    breakable,
    before={\par\Needspace{0.36\textheight}}
}

\newtcolorbox{respuestabox}{
    colback=celesteSuave,
    colframe=azulFuerte,
    boxrule=0.8pt,
    arc=4pt
}

\newcommand{\portadasolucionario}[2]{
    \thispagestyle{plain}
    \AddToShipoutPicture*{%
        \put(54, 700){\includegraphics[width=2cm]{\logoup}}
        \put(420, 706){\includegraphics[width=5cm]{\logousach}}
    }
    \vspace*{0.5cm}
    \begin{center}
        \color{azulFuerte}
        {\huge \textbf{SOLUCIONARIO}} \\[0.5cm]
        {\Huge \textbf{\curso}} \\[0.3cm]
        {\Large #1} \\[0.3cm]
        {\large #2} \\[1.4cm]
        \begin{tcolorbox}[colback=celesteSuave, colframe=azulFuerte, arc=10pt,
            width=0.82\textwidth, center, boxrule=1.5pt]
            \centering
            \Large\textbf{Desarrollo paso a paso}
        \end{tcolorbox}
    \end{center}
    \vspace{1.4cm}
    \begin{center}
        {\Large \textbf{\semestre}} \\[0.7cm]
        {\large \textbf{\autor\ -- \institucion}}
    \end{center}
    \vfill
    \begin{center}
        {\color{azulFuerte}\hrule height 0.5pt}
        \vspace{0.3cm}
        \begin{minipage}{0.45\textwidth}
            \centering
            \textbf{Instagram:} \\
            \small\href{\instagraminstitucion}{@usach.premium}
        \end{minipage}
        \begin{minipage}{0.45\textwidth}
            \centering
            \textbf{Web:} \\
            \small\href{http://\webinstitucion}{\webinstitucion}
        \end{minipage}
    \end{center}
    \newpage
    \MarcaAgua
}

\newcommand{\respuesta}[1]{
    \begin{respuestabox}
        \textbf{Respuesta:} #1
    \end{respuestabox}
}


\begin{document}
\portadasolucionario{Gu\'ia de Preparaci\'on PEP 1}{Secci\'on I -- N\'umeros reales e inecuaciones}

\begin{solucionbox}{Ejercicio 1.1}
Resuelva la inecuaci\'on
\[
    \frac{x-4}{x^2-9}\geq 0.
\]

\textbf{Soluci\'on:}

Primero factorizamos el denominador:
\[
    x^2-9=(x-3)(x+3).
\]
Debemos considerar la restricci\'on
\[
    x\neq -3,\qquad x\neq 3.
\]
Entonces la inecuaci\'on queda
\[
    \frac{x-4}{(x-3)(x+3)}\geq 0.
\]
Los puntos cr\'iticos son $x=-3$, $x=3$ y $x=4$. Construimos la tabla de signos:

\[
\begin{array}{c|cccc}
    & (-\infty,-3) & (-3,3) & (3,4) & (4,+\infty) \\ \hline
    x-4 & - & - & - & + \\
    x-3 & - & - & + & + \\
    x+3 & - & + & + & + \\ \hline
    \dfrac{x-4}{(x-3)(x+3)} & - & + & - & +
\end{array}
\]

Como buscamos que la expresi\'on sea mayor o igual que cero, tomamos los intervalos donde el signo es positivo y agregamos el punto que anula el numerador, $x=4$.

\respuesta{$x\in(-3,3)\cup[4,+\infty)$.}
\end{solucionbox}

\begin{solucionbox}{Ejercicio 1.2}
Determine el conjunto soluci\'on de
\[
    \frac{|x-1|-3}{x+2}<0.
\]

\textbf{Soluci\'on:}

Primero observamos que el denominador no puede ser cero:
\[
    x+2\neq 0 \quad\Longrightarrow\quad x\neq -2.
\]

Ahora analizamos el numerador:
\[
    |x-1|-3=0
    \quad\Longleftrightarrow\quad
    |x-1|=3.
\]
Por propiedad del valor absoluto,
\[
    x-1=3 \quad \vee \quad x-1=-3,
\]
por lo tanto
\[
    x=4 \quad \vee \quad x=-2.
\]

Los puntos cr\'iticos son $x=-2$ y $x=4$. Adem\'as, $x=-2$ no pertenece al dominio de la expresi\'on.

\[
\begin{array}{c|ccc}
    & (-\infty,-2) & (-2,4) & (4,+\infty) \\ \hline
    |x-1|-3 & + & - & + \\
    x+2 & - & + & + \\ \hline
    \dfrac{|x-1|-3}{x+2} & - & - & +
\end{array}
\]

Como la inecuaci\'on es estricta, no se incluyen los ceros del numerador. Adem\'as, $x=-2$ queda excluido por restricci\'on.

\respuesta{$x\in(-\infty,-2)\cup(-2,4)$.}
\end{solucionbox}

\begin{solucionbox}{Ejercicio 1.3}
Resuelva en $\R$:
\[
    \frac{x+2}{x-1}\leq \frac{x-3}{x+4}.
\]

\textbf{Soluci\'on:}

Primero anotamos las restricciones:
\[
    x\neq 1,\qquad x\neq -4.
\]
Llevamos todo a un mismo lado:
\[
    \frac{x+2}{x-1}\leq \frac{x-3}{x+4}
    \quad\Longleftrightarrow\quad
    0\leq \frac{x-3}{x+4}-\frac{x+2}{x-1}.
\]
Calculamos la diferencia:
\[
    0\leq
    \frac{(x-3)(x-1)-(x+2)(x+4)}{(x+4)(x-1)}.
\]
Desarrollando el numerador:
\[
    (x-3)(x-1)=x^2-4x+3,
\]
\[
    (x+2)(x+4)=x^2+6x+8.
\]
Luego,
\[
    (x-3)(x-1)-(x+2)(x+4)
    =x^2-4x+3-x^2-6x-8=-10x-5.
\]
Por tanto:
\[
    0\leq \frac{-10x-5}{(x+4)(x-1)}
    =\frac{-5(2x+1)}{(x+4)(x-1)}.
\]
Como $-5<0$, esta desigualdad es equivalente a
\[
    \frac{2x+1}{(x+4)(x-1)}\leq 0.
\]

Los puntos cr\'iticos son
\[
    x=-4,\qquad x=-\frac12,\qquad x=1.
\]

\[
\begin{array}{c|cccc}
    & (-\infty,-4) & (-4,-\frac12) & (-\frac12,1) & (1,+\infty) \\ \hline
    2x+1 & - & - & + & + \\
    x+4 & - & + & + & + \\
    x-1 & - & - & - & + \\ \hline
    \dfrac{2x+1}{(x+4)(x-1)} & - & + & - & +
\end{array}
\]

Buscamos signo menor o igual que cero. Se incluye $x=-\frac12$ porque anula el numerador y no anula el denominador.

\respuesta{$x\in(-\infty,-4)\cup\left[-\dfrac12,1\right)$.}
\end{solucionbox}

\begin{solucionbox}{Ejercicio 1.4}
El largo de un terreno rectangular supera al ancho en $7$ metros. Determine las dimensiones posibles si el \'area debe ser a lo m\'as $260\,m^2$.

\textbf{Soluci\'on:}

Designemos por $x$ al ancho del terreno, medido en metros. Como el largo supera al ancho en $7$ metros, el largo es
\[
    x+7.
\]
Como el \'area debe ser a lo m\'as $260\,m^2$, se debe cumplir
\[
    x(x+7)\leq 260.
\]
Desarrollando:
\[
    x^2+7x\leq 260
    \quad\Longleftrightarrow\quad
    x^2+7x-260\leq 0.
\]
Factorizamos:
\[
    x^2+7x-260=(x+20)(x-13).
\]
Entonces debemos resolver
\[
    (x+20)(x-13)\leq 0.
\]

Los puntos cr\'iticos son $x=-20$ y $x=13$. Como la par\'abola abre hacia arriba, la expresi\'on es menor o igual que cero entre sus ra\'ices:
\[
    x\in[-20,13].
\]
Pero $x$ representa una longitud, por lo tanto
\[
    x\geq 0.
\]
Intersectando:
\[
    x\in[-20,13]\cap[0,+\infty)=[0,13].
\]
Luego, el ancho puede variar entre $0$ y $13$ metros. El largo es $x+7$, por lo que var\'ia entre $7$ y $20$ metros.

\respuesta{Ancho entre $0$ y $13$ m; largo entre $7$ y $20$ m.}
\end{solucionbox}

\begin{solucionbox}{Ejercicio 1.5}
Resuelva la inecuaci\'on
\[
    \sqrt{x+3}\leq x-1.
\]

\textbf{Soluci\'on:}

Primero consideramos las restricciones. Para que la ra\'iz est\'e definida:
\[
    x+3\geq 0 \quad\Longrightarrow\quad x\geq -3.
\]
Adem\'as, como $\sqrt{x+3}\geq 0$, para que se cumpla
\[
    \sqrt{x+3}\leq x-1
\]
necesitamos que
\[
    x-1\geq 0
    \quad\Longrightarrow\quad
    x\geq 1.
\]
Bajo esta restricci\'on podemos elevar al cuadrado ambos lados:
\[
    \sqrt{x+3}\leq x-1
    \quad\Longleftrightarrow\quad
    x+3\leq (x-1)^2.
\]
Desarrollando:
\[
    x+3\leq x^2-2x+1
    \quad\Longleftrightarrow\quad
    0\leq x^2-3x-2.
\]
Resolvemos la ecuaci\'on asociada:
\[
    x^2-3x-2=0.
\]
Por la f\'ormula cuadr\'atica:
\[
    x=\frac{3\pm\sqrt{(-3)^2-4(1)(-2)}}{2}
     =\frac{3\pm\sqrt{17}}{2}.
\]
Como la par\'abola abre hacia arriba:
\[
    x^2-3x-2\geq 0
    \quad\Longleftrightarrow\quad
    x\in\left(-\infty,\frac{3-\sqrt{17}}{2}\right]\cup
    \left[\frac{3+\sqrt{17}}{2},+\infty\right).
\]
Finalmente, debemos intersectar con $x\geq 1$:
\[
    x\in\left[\frac{3+\sqrt{17}}{2},+\infty\right).
\]

\respuesta{$x\in\left[\dfrac{3+\sqrt{17}}{2},+\infty\right)$.}
\end{solucionbox}

\begin{solucionbox}{Ejercicio 1.6}
Determine el conjunto soluci\'on de
\[
    \sqrt{x^2-4x}>x-1.
\]

\textbf{Soluci\'on:}

Primero hallamos el dominio:
\[
    x^2-4x\geq 0
    \quad\Longleftrightarrow\quad
    x(x-4)\geq 0.
\]
De la tabla de signos se obtiene
\[
    \Dom=\left(-\infty,0\right]\cup[4,+\infty).
\]

Ahora analizamos por casos.

\textbf{Caso 1: $x\leq 0$.}

Si $x\leq 0$, entonces
\[
    x-1<0.
\]
Pero
\[
    \sqrt{x^2-4x}\geq 0.
\]
Por lo tanto,
\[
    \sqrt{x^2-4x}>x-1
\]
se cumple para todo $x\leq 0$.

\textbf{Caso 2: $x\geq 4$.}

Si $x\geq 4$, entonces $x-1>0$, y podemos elevar al cuadrado:
\[
    \sqrt{x^2-4x}>x-1
    \quad\Longleftrightarrow\quad
    x^2-4x>(x-1)^2.
\]
Desarrollando:
\[
    x^2-4x>x^2-2x+1.
\]
Reduciendo:
\[
    -4x>-2x+1
    \quad\Longleftrightarrow\quad
    -2x>1
    \quad\Longleftrightarrow\quad
    x<-\frac12.
\]
Esto contradice la condici\'on $x\geq 4$, por lo que no hay soluciones en este caso.

\respuesta{$x\in(-\infty,0]$.}
\end{solucionbox}

\end{document}
